%%=============================================================================
%% CAPITULO 1 -- INTRODUCAO
%%=============================================================================

\section{Introdução}

Modelos de linguagem de grande escala passaram a ser usados na administração
pública para classificar documentos, uma tarefa repetitiva e de alto volume em
que a automação promete ganho de escala. Esses modelos, porém, não eliminam o
erro, e um erro de classificação em um processo público pode afetar a
recuperação de informações, a consistência das decisões e a prestação de
contas. A questão prática deixa de ser se o modelo acerta em média e passa a
ser em quais casos a sua resposta pode ser aceita sem conferência.

Esta pesquisa parte da pergunta: em que medida sinais de incerteza produzidos
pelos próprios modelos de linguagem permitem decidir quando uma classificação
pode ser aceita automaticamente e quando deve ser encaminhada à revisão
humana? A pergunta aplica ao contexto da administração pública a ideia de
classificação seletiva, em que o classificador pode abster-se nos casos
duvidosos \parencite{elyaniv2010}. O objetivo geral é avaliar se esses sinais
sustentam um mecanismo de encaminhamento que reduza o erro entre os casos
aceitos sem transferir todo o trabalho de volta às pessoas.
